\subsection{CHARACTERIZATIONS OF DECOMPOSABLE LIE ALGEBRAS}

Every decomposable Lie algebra is generated as a vector space (and a fortiori as a Lie algebra) by the set of its elements that are either semi-simple or nilpotent. Conversely:

THEOREM 1. Let g be a Lie subalgebra of \(\mathfrak{gl}(V)\) and let X be a subset of g generating g as a Lie algebra over k. If every element of X is either semi-simple or nilpotent, g is decomposable.

\begin{enumerate}
\def\labelenumi{\alph{enumi})}
\tightlist
\item
  \(\mathfrak{g}\) is commutative.
\end{enumerate}

The semi-simple (resp. nilpotent) elements of g form a vector subspace \(g_{s}\) (resp. \(g_{n}\) ). The assumption is equivalent to \(g = g_{s} \oplus g_{n}\) , hence the fact that g is decomposable.

\begin{enumerate}
\def\labelenumi{\alph{enumi})}
\setcounter{enumi}{1}
\tightlist
\item
  \(\mathfrak{g}\) is reductive.
\end{enumerate}

Then \(\mathfrak{g} = \mathfrak{g}' \times \mathfrak{c}\) with \(\mathfrak{g}'\) semi-simple and \(\mathfrak{c}\) commutative. By Prop. 2, \(\mathfrak{g}'\) is decomposable. Let \(x = a + b \in \mathfrak{g}\) with \(a \in \mathfrak{g}'\) , \(b \in \mathfrak{c}\) . Let \(a_s, a_n, b_s, b_n\) be the semi-simple and nilpotent components of \(a, b\) . Since \(a_s, a_n, b_s, b_n\) mutually commute, the semi-simple and nilpotent components of \(x\) are \(a_s + b_s, a_n + b_n\) . Now \(a_s, a_n \in \mathfrak{g}'\) . If \(x\) is semi-simple, \(x = a_s + b_s\) ; since \(a_s \in \mathfrak{g}'\) , we have \(b_s \in \mathfrak{g}\) , so \(b_s \in \mathfrak{c}\) since \(b_s\) commutes with \(\mathfrak{g}\) ; consequently, \(a = a_s\) and \(b = b_s\) . Similarly, if \(x\) is nilpotent, \(a = a_n\) and \(b = b_n\) . It follows that the projections on \(\mathfrak{c}\) of the elements of X are either semi-simple or nilpotent; by \(a)\) , this implies that \(\mathfrak{c}\) is decomposable. Retaining the preceding notation, but without the assumption on \(x\) , we now have \(b_s, b_n \in \mathfrak{c}\) , so \(a_s + b_s, a_n + b_n \in \mathfrak{g}\) , which proves the theorem in this case.

\begin{enumerate}
\def\labelenumi{\alph{enumi})}
\setcounter{enumi}{2}
\tightlist
\item
  General case.
\end{enumerate}

We assume that the theorem is proved for Lie algebras of dimension \(< \dim \mathfrak{g}\) and prove it for \(\mathfrak{g}\) .

Let \(\mathfrak{n}\) be the largest ideal of nilpotency of the identity representation of \(\mathfrak{g}\) . If \(\mathfrak{n} = 0\) , \(\mathfrak{g}\) has an injective semi-simple representation, and so is reductive. Assume that \(\mathfrak{n} \neq 0\) . Let \(\mathfrak{p}\) be the normalizer of \(\mathfrak{n}\) in \(\mathfrak{gl}(\mathrm{V})\) . There exist E, \(\rho, \mathrm{F}\) satisfying the conditions of Prop. 8. Since \(\mathfrak{g} \subset \mathfrak{p}\) , \(\rho(\mathfrak{g})\) leaves F stable; let \(\rho_0\) be the representation \(u \mapsto \rho(u)|\mathrm{F}\) of \(\mathfrak{g}\) on F; we have \(\mathfrak{n} = \mathrm{Ker}\rho_0\) . The image under \(\rho\) of every semi-simple (resp. nilpotent) element of \(\mathfrak{gl}(\mathrm{V})\) is semi-simple (resp. nilpotent) (Prop. 2). The algebra \(\rho_0(\mathfrak{g})\) is thus generated by its semi-simple elements and its nilpotent elements. By the induction hypothesis, \(\rho_0(\mathfrak{g})\) is decomposable.

Let \(x \in \mathfrak{g}\) , and let \(x_s, x_n\) be its semi-simple and nilpotent components. By Prop. 2, the semi-simple and nilpotent components of \(\rho(x)\) are \(\rho(x_s), \rho(x_n)\) . Since \(\rho_0(\mathfrak{g})\) is decomposable, there exist \(y, z \in \mathfrak{g}\) such that

\[
\rho_ {0} (y) = \rho (x _ {s}) | \mathrm{F}, \rho_ {0} (z) = \rho (x _ {n}) | \mathrm{F}.
\]

Then \(x_{s}\in y + \mathfrak{n},x_{n}\in z + \mathfrak{n}\) , so \(x_{s},x_{n}\in \mathfrak{g}\)

Q.E.D.

COROLLARY 1. Every subalgebra of \(\mathfrak{gl}(\mathrm{V})\) generated by its decomposable subalgebras is decomposable.

This is clear.

COROLLARY 2. Let \(\mathfrak{g}\) be a Lie subalgebra of \(\mathfrak{gl}(\mathrm{V})\) . Then \([\mathfrak{g},\mathfrak{g}]\) is decomposable.

Let \(\mathfrak{r}\) be the radical of \(\mathfrak{g}\) , \(\mathfrak{s}\) a Levi subalgebra of \(\mathfrak{g}\) (Chap. I, §6, no. 8). Then

\[
[ \mathfrak {g}, \mathfrak {g} ] = [ \mathfrak {s}, \mathfrak {s} ] + [ \mathfrak {s}, \mathfrak {r} ] + [ \mathfrak {r}, \mathfrak {r} ] = \mathfrak {s} + [ \mathfrak {g}, \mathfrak {r} ].
\]

The algebra \([g,r]\) is decomposable since all of its elements are nilpotent (Chap. I, §5, no. 3). On the other hand, s is decomposable (Prop. 2). It follows that \([g,g]\) is decomposable (Cor. 1).

COROLLARY 3. Let \(\mathfrak{g}\) be a Lie subalgebra of \(\mathfrak{gl}(\mathrm{V})\) , and let \(\mathrm{X}\) be a subset of \(\mathfrak{g}\) generating \(\mathfrak{g}\) (as a Lie algebra over \(k\) ).

\begin{enumerate}
\def\labelenumi{(\roman{enumi})}
\item
  The decomposable envelope \(e(\mathfrak{g})\) of \(\mathfrak{g}\) is generated by the semi-simple and nilpotent components of the elements of \(X\) .
\item
  If \(k'\) is an extension of \(k\) , \(e(\mathfrak{g} \otimes_k k') = e(\mathfrak{g}) \otimes_k k'\) ; and \(\mathfrak{g}\) is decomposable if and only if \(\mathfrak{g} \otimes_k k'\) is decomposable.
\end{enumerate}

Let \(\tilde{\mathfrak{g}}\) be the subalgebra of \(\mathfrak{gl}(\mathrm{V})\) generated by the semi-simple and nilpotent components of the elements of X. Then \(\mathfrak{g} \subset \tilde{\mathfrak{g}} \subset e(\mathfrak{g})\) ; by Th. 1, \(\tilde{\mathfrak{g}}\) is decomposable, so \(\tilde{\mathfrak{g}} = e(\mathfrak{g})\) , which proves (i). Assertion (ii) follows, since X generates the \(k'\) -algebra \(\mathfrak{g} \otimes_k k'\) .

COROLLARY 4. Let g be a decomposable Lie subalgebra of \(\mathfrak{gl}(\mathrm{V})\) . Let T be the set of commutative subalgebras of g consisting of semi-simple elements (cf.~Prop. 6). The maximal elements of T all have the same dimension.

Let \(k'\) be an algebraically closed extension of \(k\) and \(V' = V \otimes_k k', \mathfrak{g}' = \mathfrak{g} \otimes_k k'\) . Let \(t_1, t_2\) be maximal elements of \(\mathcal{T}\) , \(t_i' = t_i \otimes_k k'\) , \(\mathfrak{h}_i\) the commutant of \(t_i\) in \(\mathfrak{g}\) , \(\mathfrak{h}_i' = \mathfrak{h}_i \otimes_k k'\) . Then \(\mathfrak{h}_i\) is a Cartan subalgebra of \(\mathfrak{g}\) (Prop. 6) so \(\mathfrak{h}_i'\) is a Cartan subalgebra of \(\mathfrak{g}'\) . Then \(\mathfrak{h}_i = t_i \times \mathfrak{n}_V(\mathfrak{h}_i)\) , hence \(\mathfrak{h}_i' = t_i' \times \mathfrak{n}_{V'}(\mathfrak{h}_i')\) , so that \(t_i'\) is the set of semi-simple elements of \(\mathfrak{h}_i'\) . Since \(\mathfrak{g}'\) is decomposable (Cor. 3), \(t_1'\) and \(t_2'\) are conjugate under \(\text{Aut}_e(\mathfrak{g}')\) (Prop. 6), so \(\dim t_1 = \dim t_2\) .

THEOREM 2. Let \(\mathfrak{g}\) be a Lie subalgebra of \(\mathfrak{gl}(\mathrm{V})\) . The following conditions are equivalent:

\begin{enumerate}
\def\labelenumi{(\roman{enumi})}
\item
  \(\mathfrak{g}\) is decomposable;
\item
  every Cartan subalgebra of \(\mathfrak{g}\) is decomposable;
\item
  \(\mathfrak{g}\) has a decomposable Cartan subalgebra;
\item
  the radical of \(\mathfrak{g}\) is decomposable.
\item
  \(\Longrightarrow\) (ii): This follows from Cor. 2 of Prop. 3.
\item
  \(\Longrightarrow\) (i): This follows from Cor. 1 of Th. 1, since \(\mathfrak{g}\) is generated by its Cartan subalgebras (§2, no. 3, Cor. 3 of Th. 1).
\item
  \(\Longrightarrow\) (iii): This is clear.
\item
  \(\Longrightarrow\) (ii): By Cor. 3 of Th. 1, we can assume that k is algebraically closed. The Cartan subalgebras of g are then conjugate under the elementary automorphisms of g (§3, no. 2, Th. 1); in view of Remark 1 of §3, no. 1, it follows that, if one of these is decomposable, they all are.
\item
  \(\Longrightarrow\) (iv): This follows from Cor. 2 of Prop. 4.
\item
  \(\Longrightarrow\) (i): Assume that the radical r of g is decomposable. Let s be a Levi subalgebra of g; it is decomposable (Prop. 2). Hence \(g = s + r\) is decomposable (Cor. 1 of Th. 1).
\end{enumerate}

